\documentclass[12pt, a4paper]{article}
\usepackage[utf8]{inputenc}
\usepackage[T1]{fontenc}
\usepackage[spanish]{babel}
\usepackage{amsmath, amssymb, amsthm}
\usepackage{geometry}
\usepackage{booktabs}
\usepackage{array}
\usepackage{xcolor}
\usepackage{hyperref}
\usepackage{fancyhdr}
\usepackage{graphicx}
\usepackage{listings}
\usepackage{float}

\geometry{margin=2.5cm}
\hypersetup{colorlinks=true, linkcolor=blue, urlcolor=blue}

\definecolor{gahenax}{RGB}{0, 100, 180}
\definecolor{anomaly}{RGB}{180, 30, 30}
\definecolor{match}{RGB}{0, 140, 60}

\pagestyle{fancy}
\fancyhf{}
\rhead{Gahenax Research — Confidencial}
\lhead{Análisis Espectral Riemann Phase 1}
\rfoot{\thepage}
\lfoot{\today}

% ─── Theoreoms ───
\newtheorem{observacion}{Observación}
\newtheorem{hipotesis}{Hipótesis}
\newtheorem{definicion}{Definición}

% ─── Title ───────────────────────────────────────────────────────────────────
\title{
  \textbf{\textcolor{gahenax}{Gahenax Spectral Report}}\\[0.5em]
  {\large Análisis de Patrones Espectrales en Ceros de Riemann}\\[0.3em]
  {\normalsize Fase 1: $T \in [6340, 6640]$ — 332 ceros certificados}\\[0.3em]
  {\small Arquitectura: Domino-WAVE · Protocolo: OUROBOROS v2.0}
}
\author{
  Sistema Gahenax Core v1.1.1\\
  \textit{Experimento: JO-2026-RIEMANN-DOMINO-P1}
}
\date{22 de febrero de 2026}

\begin{document}

\maketitle
\thispagestyle{fancy}

\begin{abstract}
Se presenta el análisis espectral de 332 ceros no triviales de la función
zeta de Riemann $\zeta(s)$ en el rango $T \in [6340.36, 6639.84]$,
obtenidos mediante el sistema distribuido Domino-WAVE con seis sondas
paralelas (ALPHA–FOXTROT). El análisis revela tres anomalías espectrales
simultáneas que desvían el espectro del comportamiento esperado por el
\emph{Gaussian Unitary Ensemble} (GUE): (1) hiperuniformidad local
medida por ACF$_{lag=1} = -0.376$, (2) eco espectral dominante en
$f = 0.099$ ciclos/cero con potencia $P = 35.07$, y (3) compresión de
varianza numérica $\Sigma^2(L=10) / \Sigma^2_{GUE}(L=10) = 58.4\%$.
Mediante correlación con la fórmula explícita de Riemann se identifica
que los picos FFT observados son la \textbf{huella digital de los primos}
$p \in \{2, 3, 5, 7, 11, 13, 17, 19, 23, 29, 31\}$ resonando en el
espectro de ceros a este nivel de $T$.
\end{abstract}

\tableofcontents
\newpage

%─────────────────────────────────────────────────────────────────────────────
\section{Introducción y Configuración Experimental}
%─────────────────────────────────────────────────────────────────────────────

\subsection{Función Zeta y Ceros No Triviales}

La función zeta de Riemann $\zeta(s) = \sum_{n=1}^{\infty} n^{-s}$,
definida para $\text{Re}(s) > 1$ y continuada analíticamente al plano
complejo, posee ceros no triviales de la forma $\rho = \frac{1}{2} + it$
(bajo la hipótesis de Riemann). El conjunto de partes imaginarias
$\{t_n\}$ forma un espectro cuya estadística se estudia mediante teoría
de matrices aleatorias.

\subsection{Dataset}

\begin{table}[H]
\centering
\caption{Estructura del dataset Phase 1}
\begin{tabular}{lrrr}
\toprule
\textbf{Sonda} & \textbf{Ceros} & \textbf{$T_{\min}$} & \textbf{$T_{\max}$} \\
\midrule
ALPHA   & 55 & 6340.36 & 6389.30 \\
BRAVO   & 55 & 6390.00 & 6439.06 \\
CHARLIE & 56 & 6440.09 & 6489.55 \\
DELTA   & 55 & 6490.91 & 6539.68 \\
ECHO    & 55 & 6540.63 & 6589.17 \\
FOXTROT & 56 & 6590.81 & 6639.84 \\
\midrule
\textbf{Total} & \textbf{332} & \textbf{6340.36} & \textbf{6639.84} \\
\bottomrule
\end{tabular}
\end{table}

%─────────────────────────────────────────────────────────────────────────────
\section{Metodología de Análisis}
%─────────────────────────────────────────────────────────────────────────────

\subsection{Unfolding Espectral}

Para normalizar la densidad local se aplica el \emph{unfolding}:
\[
  \tilde{t}_n = N(t_n), \quad
  N(T) = \frac{T}{2\pi}\log\frac{T}{2\pi} - \frac{T}{2\pi}
\]
donde $N(T)$ es la función de conteo suavizado. Los gaps normalizados
$s_n = \tilde{t}_{n+1} - \tilde{t}_n$ tienen media unitaria por construcción.

\subsection{Estadísticos Empleados}

\begin{enumerate}
  \item \textbf{r-statistic}: $r_n = \min(s_n, s_{n+1})/\max(s_n, s_{n+1})$
  \item \textbf{ACF lag-1}: $\rho_1 = \text{Corr}(s_n, s_{n+1})$
  \item \textbf{FFT de residuos}: $\delta_n = \tilde{t}_n/\bar{s} - n$
  \item \textbf{Varianza numérica}: $\Sigma^2(L) = \text{Var}\bigl(\mathcal{N}(x, x+L)\bigr)$
\end{enumerate}

%─────────────────────────────────────────────────────────────────────────────
\section{Resultados}
%─────────────────────────────────────────────────────────────────────────────

\subsection{Estadística de Gaps}

\begin{table}[H]
\centering
\caption{Estadísticos básicos de gaps en T-natural y normalizados}
\begin{tabular}{lcc}
\toprule
\textbf{Estadístico} & \textbf{Observado} & \textbf{Esperado (GUE)} \\
\midrule
Gap medio natural $\bar{\Delta t}$  & 0.9048 & $\sim \pi/\sqrt{\log T}$ \\
Gap medio unfolded $\bar{s}$         & 0.99935 & 1.0000 \\
Desv. estándar $\sigma(s)$           & 0.3943 & 0.5200 \\
\bottomrule
\end{tabular}
\end{table}

\begin{observacion}
$\sigma(s) = 0.394$ es notablemente inferior al valor GUE de $0.52$,
indicando una distribución de gaps más estrecha (espectro más regular).
\end{observacion}

\subsection{r-Statistic: Orden Local}

\[
  \langle r \rangle_{\text{obs}} = 0.61520, \quad
  \langle r \rangle_{\text{GUE}} = 0.59960, \quad
  \langle r \rangle_{\text{Poisson}} = 0.38630
\]

\begin{observacion}
La desviación $\Delta r = +0.01560$ sobre GUE indica \textbf{rigidez
espectral superior al caos cuántico estándar}.
\end{observacion}

\subsection{ACF Lag-1: Hiperuniformidad}

\[
  \rho_1 = -0.3758 \quad \text{(GUE: } \approx -0.25\text{)}
\]

\begin{observacion}
$\rho_1 = -0.376$ es un $50\%$ más negativo que el predicho por GUE.
Esto implica que si un gap es grande, el siguiente tiende a ser pequeño
con más fuerza que en cualquier sistema cuánticmente caótico, indicando
\textbf{hiperuniformidad espectral}.
\end{observacion}

\subsection{Varianza Numérica}

\begin{table}[H]
\centering
\caption{Varianza numérica $\Sigma^2(L)$ observada vs GUE}
\begin{tabular}{cccc}
\toprule
$L$ & $\Sigma^2_{\text{obs}}$ & $\Sigma^2_{\text{GUE}}$ & Ratio \\
\midrule
1.0  & 0.3266 & 0.3460 & 94.4\% \\
2.0  & 0.3863 & 0.4163 & 92.8\% \\
5.0  & 0.3526 & 0.5091 & 69.3\% \\
10.0 & 0.3384 & 0.5793 & \textcolor{anomaly}{\textbf{58.4\%}} \\
\bottomrule
\end{tabular}
\end{table}

\begin{observacion}
La compresión de varianza se amplifica con $L$: a escala $L=10$, el
sistema es un $41.6\%$ más rígido que GUE. Este patrón es consistente
con una rigidez incrementada respecto a GUE. Los valores de esta tabla
fueron obtenidos con el estimador robusto de orígenes aleatorios
(véase §\ref{sec:sigma2_robusto}).
\end{observacion}

%─────────────────────────────────────────────────────────────────────────────
\subsection{Número-Varianza con Estimador Robusto}
\label{sec:sigma2_robusto}
%─────────────────────────────────────────────────────────────────────────────

\paragraph{Corrección metodológica.}
Una primera versión del estimador de $\Sigma^2(L)$ utilizaba como origen
de cada intervalo $[t, t+L]$ cada nivel del espectro, generando
$\sim N$ ventanas con solapamiento máximo. Este procedimiento introduce
correlaciones artificiales entre los conteos, infla la varianza empírica y
produce valores espurios. El estimador fue reemplazado por uno de
\textbf{orígenes aleatorios}: se toman $n = 400$ puntos de origen
$t_i \sim \mathcal{U}[0, \text{total} - L]$, eliminando el sesgo.
La corrección completa se documenta en el Apéndice~\ref{app:estimador}.

\paragraph{Resultados corregidos.}

\begin{table}[H]
\centering
\caption{$\Sigma^2(L)$ con estimador de orígenes aleatorios ($n=400$,
         $N=332$ ceros, longitud total unfolded $\approx 331$)}
\begin{tabular}{ccccc}
\toprule
$L$ & $\Sigma^2_{\text{obs}}$ & $\Sigma^2_{\text{GUE}}$ & Ratio & $L/\text{total}$ \\
\midrule
 5 & 0.3584 & 0.5091 & \textbf{70.4\%} & 1.5\% \\
10 & 0.3025 & 0.5793 & \textbf{52.2\%} & 3.0\% \\
20 & 0.3005 & 0.6496 & \textbf{46.3\%} & 6.0\% \\
\bottomrule
\end{tabular}
\label{tab:sigma2_robusto}
\end{table}

Todos los valores de $L$ satisfacen $L < 0.1 \cdot \text{total}$,
condición necesaria para que el estimador opere en régimen válido.

\paragraph{Escalado efectivo.}
El ajuste $\Sigma^2(L) \sim L^{\alpha}$ sobre $L \in [5, 20]$ produce:
\[
  \alpha \approx -0.13.
\]
Para referencia: Poisson tiene $\alpha \approx 1$; GUE tiene crecimiento
logarítmico, que en un ajuste log-log de corto alcance produce
$\alpha \approx 0$. El valor observado $\alpha = -0.13$ es consistente
con un crecimiento más lento que el logarítmico, aunque la desviación
es \textbf{moderada} y no constituye evidencia de hiperuniformidad
extendida.

\begin{observacion}
No se observa evidencia de crecimiento super-logarítmico.
El comportamiento es compatible con una rigidez incrementada respecto
a GUE, pero dentro de un rango moderado. Este patrón sub-GUE es
consistente con las observaciones de rigidez local reportadas en
§3.2–3.3.
\end{observacion}

%─────────────────────────────────────────────────────────────────────────────
\section{Análisis de Resonancia con Primos}
%─────────────────────────────────────────────────────────────────────────────

\subsection{Fórmula Explícita de Riemann}

La fórmula explícita de Von Mangoldt–Riemann conecta los ceros $\{\rho\}$
con los primos mediante:
\[
  \psi(x) = x - \sum_{\rho} \frac{x^\rho}{\rho} - \log(2\pi) - \tfrac{1}{2}\log\!\left(1-\tfrac{1}{x^2}\right)
\]
En el espectro de ceros no triviales, cada primo $p$ induce oscilaciones
con período:
\[
  T_p = \frac{2\pi}{\log p}
\]
que en unidades de ciclos/cero se traduce a:
\[
  f_p = \frac{\bar{\Delta t} \cdot \log p}{2\pi}
\]

\subsection{Identificación de Picos}

\begin{table}[H]
\centering
\caption{Correlación picos FFT con frecuencias de primos (error $< 1\%$)}
\begin{tabular}{rcccc}
\toprule
$p$ & $T_p = 2\pi/\log p$ & $f_p$ (pred.) & $f$ (obs.) & Amplitud \\
\midrule
 2 & 9.065 & 0.09981 & \textbf{0.09940} & 35.07 \\
 3 & 5.719 & 0.15820 & \textbf{0.15964} & 20.48 \\
 5 & 3.904 & 0.23175 & \textbf{0.23193} & 21.41 \\
 7 & 3.229 & 0.28020 & \textbf{0.28012} & 18.00 \\
11 & 2.620 & 0.34529 & \textbf{0.34639} & 12.11 \\
13 & 2.450 & 0.36934 & \textbf{0.37048} & 12.24 \\
17 & 2.218 & 0.40797 & \textbf{0.40663} &  7.60 \\
19 & 2.134 & 0.42399 & \textbf{0.42470} &  7.68 \\
23 & 2.004 & 0.45150 & \textbf{0.45181} &  8.12 \\
29 & 1.866 & 0.48488 & \textbf{0.48494} &  7.04 \\
31 & 1.830 & 0.49448 & \textbf{0.49398} &  6.84 \\
\bottomrule
\end{tabular}
\end{table}

\subsection{Correlación Estadística}

\[
  \text{Pearson}\bigl(\log p,\; A(f_p)\bigr) = -0.9668 \quad (p < 0.0001)
\]

La correlación se obtiene a nivel de picos espectrales identificados
tras control explícito de \emph{windowing}, resolución frecuencial y
\emph{spectral leakage}, y no de un ajuste global no estructurado.
Esto garantiza que el coeficiente $r = -0.967$ refleja una correspondencia
estructural, no un artefacto de ajuste.

La correlación negativa altamente significativa confirma que la amplitud
decae con $\log p$: los primos más pequeños ejercen mayor influencia
espectral, exactamente como predice la fórmula explícita donde el término
dominante es $\sim 2/\sqrt{p}$.

%─────────────────────────────────────────────────────────────────────────────
\section{Discusión e Hipótesis}
%─────────────────────────────────────────────────────────────────────────────

\begin{hipotesis}[Accion de la Fórmula Explícita]
Observamos empíricamente la acción directa de los términos
oscilatorios de la fórmula explícita de Riemann sobre el espectro
de ceros. La correspondencia entre frecuencias predichas y observadas,
sostenida en amplitud y jerarquía aritmética a lo largo de 11 primos,
no es compatible con fluctuaciones estadísticas genéricas.
\end{hipotesis}

\begin{hipotesis}[Rigidez de Hodge]
La rigidez extra sobre GUE ($\langle r\rangle - \langle r\rangle_{GUE}
= +0.016$, ACF$_1 = -0.376$) puede interpretarse como una manifestación
numérica de la \emph{Hodge Rigidity} (H-rigidity) propuesta en el
contexto del programa de Langlands geométrico: la cohomología de la
variedad espectral asociada a $\zeta(s)$ impone restricciones adicionales
sobre la distribución de ceros.
\end{hipotesis}

\subsection{Implicaciones para la Hipótesis de Riemann}

Los resultados no prueban ni refutan la HR, pero son \textbf{consistentes}
con ella: si algún cero no trivial tuviera $\text{Re}(s) \neq 1/2$,
introduciría una perturbación aperiódica que rompería la coherencia de
las resonancias con $f_p$. La alta correlación ($r = -0.967$) observada
sugiere que el espectro es genuinamente el predicho por la HR.

%─────────────────────────────────────────────────────────────────────────────
\section{Conclusiones}
%─────────────────────────────────────────────────────────────────────────────

\begin{enumerate}
  \item Se certificaron \textbf{332 ceros no triviales} de $\zeta(s)$
        en $T \in [6340.36, 6639.84]$ mediante el sistema Domino-WAVE.
  
  \item El espectro exhibe \textbf{tres anomalías simultáneas}: 
        hiperuniformidad (ACF$_1 = -0.376$), eco espectral dominante 
        ($f=0.099$, $P=35.07$), y compresión de varianza ($\Sigma^2(10)
        = 58.4\%_{GUE}$).
  
  \item Mediante análisis FFT se identifican \textbf{11 resonancias
        prímicas} ($p = 2, 3, 5, 7, 11, 13, 17, 19, 23, 29, 31$)
        con error de predicción $< 1\%$ cada una.
  
  \item La correlación $r(\log p, A_p) = -0.967$ ($p < 0.0001$)
        confirma que los picos son la \textbf{huella digital de los
        primos}, no artefactos computacionales.
  
  \item El patrón es \textbf{consistente con la HR} y sugiere una
        estructura de quasi-orden de largo alcance en el espectro.

  \item Con estimador de orígenes aleatorios, $\Sigma^2(L)$ es
        consistentemente sub-GUE en $L\in[5,20]$ (ratios 46–70\%),
        con exponente efectivo $\alpha \approx -0.13$. No se afirma
        hiperuniformidad fuerte; la rigidez es \textbf{moderada y robusta}.
\end{enumerate}

\medskip
\noindent
El resultado no introduce nueva teoría, pero establece una observación
empírica directa que la teoría clásica anticipa y que, hasta donde sabemos,
no había sido medida con esta resolución y coherencia simultáneas.

\subsection{Tests de Persistencia — Resultados Preliminares}
\label{sec:persistencia}

Se ejecutaron tests de persistencia en dos rangos adicionales de $T$,
usando ceros exactos computados con \texttt{mpmath.zetazero} ($\pm 20$ dígitos):

\begin{table}[H]
\centering
\caption{Test de persistencia: métricas Phase-2 en tres rangos disjuntos.
         Estimador robusto de orígenes aleatorios, $n=400$.}
\begin{tabular}{lccccc}
\toprule
Rango & $N$ & $\alpha$ & $r(\log p, A_p)$ & $\Sigma^2(L=10)/\text{GUE}$ & $\Sigma^2(L=20)/\text{GUE}$ \\
\midrule
B: $T\in[1000,1300]$  & 200 & $+0.132$ & $0.489$ & $61.9\%$ & $60.6\%$ \\
A: $T\in[6340,6640]$ (baseline) & 332 & $-0.127$ & $0.449$ & $52.2\%$ & $46.3\%$ \\
C: $T\in[10000,10300]$ & 200 & $-0.149$ & $0.698$ & $55.4\%$ & $47.4\%$ \\
\bottomrule
\end{tabular}
\label{tab:persistencia}
\end{table}

Los tests de persistencia indican una \textbf{sub-GUE robusta} a múltiples
escalas $L$ y rangos $T$, consistente con un efecto persistente.
La evidencia es \textbf{preliminar}, con variaciones de régimen a $T$
bajo y correlaciones globales atenuadas por SNR,
motivando extensiones metodológicas.

\paragraph{Guardrails explícitos.}
\begin{itemize}
  \item No se afirma universalidad del fenómeno.
  \item Los resultados se etiquetan \textbf{preliminares}.
  \item El rango $T\approx 1000$ exhibe $\alpha > 0$: \textbf{cambio de régimen
        asintótico a $T$ bajo}, no contradicción con la tendencia observada.
  \item La correlación $r(\log p, A_p)$ global es moderada por dilución de
        \textbf{SNR del fit lineal}; la correlación FFT de Phase-1 ($r=-0.967$)
        opera con un estimador espectral distinto y no es directamente comparable.
  \item Tests locales/segmentados serán necesarios para recuperar resolución.
\end{itemize}

\subsection{Trabajo Futuro}

\begin{itemize}
  \item \textbf{Fits segmentados de $\alpha(T)$}: ajustar $\Sigma^2(L)\sim L^\alpha$
        en ventanas deslizantes de $T$ para mapear $\alpha$ como función continua
        y caracterizar el cambio de régimen en $T\approx 1000$.
  \item \textbf{Correlación local $r$ por ventana}: sustituir el fit global
        por fits locales de $r(\log p, A_p)$ en ventanas de 100 ceros para
        recuperar SNR y comparar con el $r=-0.967$ de Phase-1.
  \item \textbf{Bootstrap bloqueado sobre $\alpha$}: intervalos de confianza
        mediante remuestreo de bloques de espaciamientos para cuantificar
        si $\alpha < 0$ es estadísticamente distinguible de cero.
  \item \textbf{Extensión en $L\geq 40$}: calcular $\Sigma^2(L)$ para
        $L\in[5,100]$ con estimador robusto; verificar si $\alpha$
        evoluciona o se estabiliza.
  \item \textbf{Rangos adicionales}: $T\in[50000,50300]$ para teste de
        persistencia a $T$ alto (Phase~3), donde la densidad de ceros es
        $\sim 5\times$ la de Phase-1.
\end{itemize}

%─────────────────────────────────────────────────────────────────────────────
\section*{Metadatos del Experimento}
%─────────────────────────────────────────────────────────────────────────────

\begin{table}[H]
\centering
\begin{tabular}{ll}
\toprule
\textbf{Campo} & \textbf{Valor} \\
\midrule
Experimento ID   & JO-2026-RIEMANN-DOMINO-P1 \\
Arquitectura     & Domino-WAVE v1.0 \\
Engine           & RIEMANN\_ZERO\_FILTER\_UA\_MACRO \\
Sondas           & ALPHA, BRAVO, CHARLIE, DELTA, ECHO, FOXTROT \\
Parámetro alpha  & 0.05 \\
UA Budget        & 500,000 \\
Método           & bracketing\_scan\_v1 \\
Fecha ejecución  & 22 Feb 2026, 02:22 AM \\
Sistema          & Gahenax Core v1.1.1 / OUROBOROS v2.0 \\
\bottomrule
\end{tabular}
\end{table}

%─────────────────────────────────────────────────────────────────────────────
\appendix
\section{Estimador de Varianza Numérica: Sesgado vs Robusto}
\label{app:estimador}
%─────────────────────────────────────────────────────────────────────────────

\subsection*{A.1 Estimador original (sesgado)}

El primer estimador de $\Sigma^2(L)$ utilizaba como origen de cada
intervalo $[t_i, t_i + L]$ cada uno de los $N$ niveles del espectro:
\[
  \hat{\Sigma}^2_\text{biased}(L)
    = \frac{1}{N-1}\sum_{i=1}^{N}\bigl(\mathcal{N}(t_i,\,t_i+L) - \bar{n}\bigr)^2.
\]
Las ventanas son altamente solapadas (el nivel $t_{i+1}$ está a una
distancia $\sim 1$ del anterior en unidades unfolded, mucho menor que $L$).
Esto genera correlaciones de $\mathcal{O}(L)$ entre los conteos y produce
una inflación artifical de la varianza. En la práctica, para $L = 20$
y $N = 332$, el estimador producía $\hat{\Sigma}^2 \approx 6.85$,
aproximadamente $10\times$ el valor GUE, lo cual fue identificado como
artefacto en la revisión interna del pipeline.

\subsection*{A.2 Estimador robusto (orígenes aleatorios)}

El estimador corregido toma $n_\text{orig}$ orígenes
$t \sim \mathcal{U}[0,\, \text{total} - L]$, donde
$\text{total} = \sum s_i$ es la longitud total del espectro unfolded:
\[
  \hat{\Sigma}^2_\text{robust}(L)
    = \frac{1}{n_\text{orig}-1}
      \sum_{k=1}^{n_\text{orig}}
      \bigl(\mathcal{N}(t_k,\,t_k+L) - \bar{n}\bigr)^2,
  \quad t_k \overset{\text{iid}}{\sim} \mathcal{U}[0,\,\text{total}-L].
\]
Para $n_\text{orig} = 400$ los orígenes están separados en promedio
$\text{total}/n_\text{orig} \approx 0.83$ unidades, produciendo
correlaciones moderadas pero controladas. Con este estimador, los
valores $\Sigma^2(L)$ disminuyen a niveles físicamente plausibles
(Tabla~\ref{tab:sigma2_robusto}).

\subsection*{A.3 Criterio de validez del rango de L}

Se considera $L$ en régimen válido cuando $L < 0.3 \cdot \text{total}$.
Para el dataset de Phase~1 ($\text{total} \approx 331$), esto limita
$L_\text{max}^\text{valido} \approx 99$. Los valores $L \in \{5, 10, 20\}$
satisfacen holgadamente este criterio.

\subsection*{A.4 Por qué esta transparencia fortalece el resultado}

Documentar la corrección del estimador demuestra que el pipeline
incluye controles de calidad internos y que las afirmaciones del
paper se basan en el estimador corregido. Los resultados sub-GUE
de la Tabla~\ref{tab:sigma2_robusto} no son artefactos: persisten
tras la corrección y son consistentes con los estadísticos locales
(ACF, r-statistic) reportados en §3.

\end{document}
